
\documentclass{article}
\usepackage{colortbl}
\usepackage{makecell}
\usepackage{multirow}
\usepackage{supertabular}

\begin{document}

\newcounter{utterance}

\centering \large Interaction Transcript for game `cladder', experiment `full\_v1.5\_default', episode 6482 with Qwen3.5{-}9B.
\vspace{24pt}

{ \footnotesize  \setcounter{utterance}{1}
\setlength{\tabcolsep}{0pt}
\begin{supertabular}{c@{$\;$}|p{.15\linewidth}@{}p{.15\linewidth}p{.15\linewidth}p{.15\linewidth}p{.15\linewidth}p{.15\linewidth}}
    \# & \multicolumn{2}{c}{Player} && \multicolumn{2}{c}{Game Master} \\
    \hline

    \theutterance \stepcounter{utterance}  
    & & & \multicolumn{4}{p{0.6\linewidth}}{
        \cellcolor[rgb]{0.9,0.9,0.9}{
            \makecell[{{p{\linewidth}}}]{
                \texttt{\tiny{[P1$\langle$GM]}}
                \texttt{You are an expert in causal inference. The following question is not a typical commonsense query, but rather a meticulously designed question created by a professor specializing in causal inference, intended to assess the students' mastery of the course content.} \\
\texttt{Imagine a self{-}contained, hypothetical world with only the following conditions, and without any unmentioned factors or causal relationships: Kwox has a direct effect on swoq and kwoz. Swoq has a direct effect on kwoz. For those who are not kwox, the probability of kwoz is 36\%. For those who are kwox, the probability of kwoz is 56\%. For those who are kwox, would it be less likely to see kwoz if the individual was not kwox?} \\
\texttt{Start your answer with 'Yes' or 'No', followed by additional reasoning or evidence to support your explanation.} \\
            }
        }
    }
     \\ \\

    \theutterance \stepcounter{utterance}  
    & \multicolumn{4}{p{0.6\linewidth}}{
        \cellcolor[rgb]{0.561,0.737,0.561}{
            \makecell[{{p{\linewidth}}}]{
                \texttt{\tiny{[P1$\rangle$GM]}}
                \texttt{No} \\
\\ 
\texttt{**Reasoning:**} \\
\\ 
\texttt{To answer this question, we must carefully parse the logical structure of the query within the constraints of the hypothetical world provided.} \\
\\ 
\texttt{1.  **Analyze the Premise:** The question asks: "For those who are **kwox**, would it be less likely to see **kwoz** if the individual was **not kwox**?"} \\
\texttt{    *   This is a counterfactual query regarding a specific group: individuals who *are* kwox.} \\
\texttt{    *   We are comparing the probability of observing `kwoz` in the actual world (where they are kwox) versus a counterfactual world (where they are not kwox).} \\
\\ 
\texttt{2.  **Analyze the Data:**} \\
\texttt{    *   \$P(\textbackslash{}text\{kwoz\} | \textbackslash{}text\{kwox\}) = 56\textbackslash{}\%\$} \\
\texttt{    *   \$P(\textbackslash{}text\{kwoz\} | \textbackslash{}neg\textbackslash{}text\{kwox\}) = 36\textbackslash{}\%\$} \\
\\ 
\texttt{3.  **Evaluate the Counterfactual:**} \\
\texttt{    *   For an individual who *is* kwox, the probability of seeing kwoz is 56\%.} \\
\texttt{    *   If that same individual were *not* kwox, the probability of seeing kwoz would be 36\% (based on the definition of the world's conditions for non{-}kwox individuals).} \\
\\ 
\texttt{4.  **Compare the Probabilities:**} \\
\texttt{    *   The question asks if it would be **less likely** to see kwoz in the counterfactual scenario (\$\textbackslash{}neg\textbackslash{}text\{kwox\}\$) compared to the actual scenario (\$\textbackslash{}text\{kwox\}\$).} \\
\texttt{    *   Is \$36\textbackslash{}\% < 56\textbackslash{}\%\$? **Yes.**} \\
\\ 
\texttt{5.  **Address the Potential Confusion (The "Trick"):**} \\
\texttt{    *   The phrasing "For those who are not kwox, the probability of kwoz is 36\%" establishes the baseline for the counterfactual state.} \\
\texttt{    *   The phrasing "For those who are kwox, the probability of kwoz is 56\%" establishes the factual state.} \\
\texttt{    *   Since 36\% is indeed less than 56\%, the event `kwoz` is less likely in the counterfactual world where the individual is not kwox.} \\
\\ 
\texttt{**Conclusion:**} \\
\texttt{Yes, for those who are kwox, it would be less likely to see kwoz if the individual was not kwox (36\% vs 56\%). The causal structure (Kwox \$\textbackslash{}to\$ Swoq \$\textbackslash{}to\$ Kwoz and Kwox \$\textbackslash{}to\$ Kwoz) explains *why* the probability is higher for kwox, but the numerical comparison directly answers the specific question asked.} \\
            }
        }
    }
    & & \\ \\

    \theutterance \stepcounter{utterance}  
    & & & \multicolumn{2}{p{0.3\linewidth}}{
        \cellcolor[rgb]{0.9,0.9,0.9}{
            \makecell[{{p{\linewidth}}}]{
                \texttt{\tiny{[GM$|$GM]}}
                \texttt{No} \\
            }
        }
    }
    & & \\ \\

    \theutterance \stepcounter{utterance}  
    & & & \multicolumn{2}{p{0.3\linewidth}}{
        \cellcolor[rgb]{0.9,0.9,0.9}{
            \makecell[{{p{\linewidth}}}]{
                \texttt{\tiny{[GM$|$GM]}}
                \texttt{yes} \\
            }
        }
    }
    & & \\ \\

    \theutterance \stepcounter{utterance}  
    & & & \multicolumn{2}{p{0.3\linewidth}}{
        \cellcolor[rgb]{0.9,0.9,0.9}{
            \makecell[{{p{\linewidth}}}]{
                \texttt{\tiny{[GM$|$GM]}}
                \texttt{game\_result = LOSE} \\
            }
        }
    }
    & & \\ \\

\end{supertabular}
}

\end{document}
